\documentclass[UTF8]{ctexart}
\usepackage{xeCJK}
\usepackage{geometry}
\geometry{left=2cm,right=2cm,top=2.5cm,bottom=2.5cm}
\setlength{\headheight}{12.7pt}

\usepackage{titlesec}
\titleformat{\section}
  {\normalfont\large\bfseries}
  {\thesection}
  {1em}
  {}
\titlespacing*{\section}
  {0pt}
  {3.5ex plus 1ex minus .2ex}
  {2.3ex plus .2ex}

\usepackage{amsmath}
\usepackage{amssymb}
\usepackage{amsthm}
\usepackage{enumitem}
\setlist[enumerate]{itemsep=0pt,topsep=0pt,parsep=0pt,leftmargin=0pt,itemindent=2.5em}
\setlist[itemize]{itemsep=0pt,topsep=0pt,parsep=0pt,leftmargin=0pt,itemindent=2.5em}
\usepackage{fancyhdr}
\pagestyle{fancy}
\fancyhf{}
\usepackage{graphicx}
\usepackage{hyperref}
\usepackage{indentfirst}
\usepackage{lmodern}


\begin{document}
\setlength{\abovedisplayskip}{1pt}
\setlength{\belowdisplayskip}{1pt}

\section{单点替换}

\hypertarget{sec:定义1.1}{}
\noindent\textbf{定义1.1 (单点替换)}

给定一阶语言$\mathcal{L}$, 变元或常元$y$和项$s$,
把\href{04 一般替换#nameddest=sec:定义1.1}{替换}
$\displaystyle i\frac{s}{y}$简记为$\displaystyle\frac{s}{y}$, 即
\[\forall z\in\mathcal{V}\cup\mathcal{C}:\quad\frac{s}{y}(z)=
\left\{\begin{aligned}
    &s,\quad&z=y,\\[-3pt]
    &z,\quad&z\neq y,
\end{aligned}\right.\]
称之为$y$到$s$的\textbf{单点替换}.

\vspace{10pt}

\hypertarget{sec:定理1.1}{}
\noindent\textbf{定理1.1}

给定一阶语言$\mathcal{L}$, 变元$x$, 变元或常元$y$, 项$s$和公式$\varphi$.
\vspace{3pt}
\begin{enumerate}
    \item 如果$y=x$, 那么$\displaystyle
    \frac{s}{y}(\exists x\varphi)=\exists x\varphi$,
    \vspace{6pt}
    \item 如果$y\neq x$, 那么$\displaystyle
    \frac{s}{y}(\exists x\varphi)=\exists x\frac{s}{y}(\varphi)$.
\end{enumerate}

\noindent{\kaishu 证明.}

依定义, $\displaystyle\frac{s}{y}(\exists x\varphi)=
\exists x\left(\frac{s}{y}/x\right)(\varphi)$.

\vspace{3pt}
\begin{enumerate}
    \item 假设$y=x$, 那么$\displaystyle\frac{s}{y}/x=i$,
    从而$\displaystyle\frac{s}{y}(\exists x\varphi)=\exists x\varphi$,
    \vspace{6pt}
    \item 假设$y\neq x$, 那么$\displaystyle\frac{s}{y}/x=\frac{s}{y}$,
    从而$\displaystyle\frac{s}{y}(\exists x\varphi)=\exists x\frac{s}{y}(\varphi)$.
    \hfill $\qedsymbol$
\end{enumerate}

\vspace{10pt}

\hypertarget{sec:定理1.2}{}
\noindent\textbf{定理1.2}

给定一阶语言$\mathcal{L}$, 变元$x$, 变元或常元$y$, 项$s$和公式$\varphi$.
\vspace{3pt}
\begin{enumerate}
    \item 如果$y=x$, 那么$\displaystyle\frac{s}{y}[\exists x\varphi]$为真,
    \vspace{6pt}
    \item 如果$y\neq x$, 那么$\displaystyle\frac{s}{y}[\exists x\varphi]$
    当且仅当$\displaystyle\frac{s}{y}[\varphi]$并且
    $y\in\mathrm{at}(\varphi)\to x\notin\mathrm{var}(s)$.
\end{enumerate}

\noindent{\kaishu 证明.}

\begin{enumerate}
    \item 假设$y=x$, 即$\displaystyle\frac{s}{y}/x=i$. 那么首先显然有$i[\varphi]$;
    其次对任意的$z\in\mathrm{at}(\exists x\varphi)$, 一定有$z\neq x$, 所以
    \[
        \frac{s}{y}(z)=z\implies\mathrm{var}\,\frac{s}{y}(z)=\{z\}
        \implies x\notin\mathrm{var}\,\frac{s}{y}(z).
    \]
    综上即得$\displaystyle\frac{s}{y}[\exists x\varphi]$.

    \vspace{3pt}
    \item 假设$y\neq x$, 即$\displaystyle\frac{s}{y}/x=\frac{s}{y}$.
    
    \vspace{6pt}\hspace{2.17em}
    如果$\displaystyle\frac{s}{y}[\exists x\varphi]$, 依定义有
    $\displaystyle\frac{s}{y}[\varphi]$并且对任意的$z\in\mathrm{at}
    (\exists x\varphi)$有$\displaystyle x\notin\mathrm{var}\frac{s}{y}(z)$. 此时
    如果$y\in\mathrm{at}(\varphi)$就有
    \[
        y\in\mathrm{at}(\exists x\varphi)\quad\implies\quad
        x\notin\mathrm{var}\frac{s}{y}(y)=\mathrm{var}(s).
    \]
    
    \hspace{2.17em}
    如果$\displaystyle\frac{s}{y}[\varphi]$并且
    $y\in\mathrm{at}(\varphi)\to x\notin\mathrm{var}(s)$, 那么
    对任意的$z\in\mathrm{at}(\exists x\varphi)$, 同样一定有$z\neq x$, 所以
    \[
        \frac{s}{y}(z)=z\implies\mathrm{var}\,\frac{s}{y}(z)=\{z\}
        \implies x\notin\mathrm{var}\,\frac{s}{y}(z).
    \]
    综上即得$\displaystyle\frac{s}{y}[\exists x\varphi]$. \hfill $\qedsymbol$
\end{enumerate}





\newpage

下面的两个定理可以简单地判断一些单点替换是自由的.

% 下面的两个定理在想象中是显然的, 但是好像并不是美妙的结论. 应该只适合作为引理使用.

% 晚上23:51, 尝试睡觉前写完后面的三个(四个?)定理.
% 2026.02.28

\vspace{10pt}

\hypertarget{sec:定理1.3}{}
\noindent\textbf{定理1.3}

给定一阶语言$\mathcal{L}$, 变元或常元$y$, 项$s$和公式$\varphi$.
如果$y\notin\mathrm{at}(\varphi)$, 那么$\displaystyle\frac{s}{y}[\varphi]$.

\noindent{\kaishu 证明.}

\begin{enumerate}
    \item 对任意的原子公式$\varphi$, 一定有$\displaystyle\frac{s}{y}[\varphi]$.
    
    \vspace{3pt}
    \item 假设$\varphi_1$满足命题, 再假设$\varphi=\neg\varphi_1$满足
    $y\notin\mathrm{at}(\varphi)$.
    那么由$\mathrm{at}(\varphi)=\mathrm{at}(\varphi_1)$得
    \vspace{3pt}
    \[
        y\notin\mathrm{at}(\varphi_1)\implies
        \frac{s}{y}[\varphi_1]\implies\frac{s}{y}[\varphi].
    \]

    \vspace{3pt}
    \item 假设$\varphi_1,\varphi_2$满足命题,
    再假设$\varphi=(\varphi_1\to\varphi_2)$满足$y\notin\mathrm{at}(\varphi)$. 那么
    由$\mathrm{at}(\varphi)=\mathrm{at}(\varphi_1)\cup\mathrm{at}(\varphi_2)$得
    \vspace{3pt}
    \[
        y\notin\mathrm{at}(\varphi_1)\land y\notin\mathrm{at}(\varphi_2)\implies
        \frac{s}{y}[\varphi_1]\land\frac{s}{y}[\varphi_2]\implies
        \frac{s}{y}[\varphi].
    \]

    \item 对任意的变元$x$, 假设$\varphi_1$满足命题,
    再假设$\varphi=\exists x\varphi_1$满足$y\notin\mathrm{at}(\varphi)$.

    \vspace{3pt}\hspace{2.17em}
    如果$y=x$, 那么$\displaystyle\frac{s}{y}[\varphi]$为真;

    \vspace{3pt}\hspace{2.17em}
    如果$y\neq x$, 那么有$y\notin\mathrm{at}(\varphi_1)$,
    依假设有$\displaystyle\frac{s}{y}[\varphi_1]$,
    直接可得$\displaystyle\frac{s}{y}[\varphi]$. \hfill $\qedsymbol$
\end{enumerate}

\vspace{10pt}

% 熬不住了, 先睡了. 刚才处理了一下编译文件位置的问题, 感谢 Copilot.
% 剩下两个定理明天再写.

% 刚刚看了《才二十三》的 MV. Good night, Khalil.
% 2026.03.01

\hypertarget{sec:定理1.4}{}
\noindent\textbf{定理1.4}

给定一阶语言$\mathcal{L}$, 变元或常元$y$, 项$s$和公式$\varphi$.
如果$\mathrm{var}(s)\cap\mathrm{bnd}(\varphi)=\varnothing$,
那么$\displaystyle\frac{s}{y}[\varphi]$.

\noindent{\kaishu 证明.}

\begin{enumerate}
    \item 对任意的原子公式$\varphi$, 一定有$\displaystyle\frac{s}{y}[\varphi]$.
    
    \vspace{3pt}
    \item 假设$\varphi_1$满足命题, 再假设$\varphi=\neg\varphi_1$满足
    $\mathrm{var}(s)\cap\mathrm{bnd}(\varphi)=\varnothing$.
    那么由$\mathrm{bnd}(\varphi)=\mathrm{bnd}(\varphi_1)$得
    \vspace{3pt}
    \[
        \mathrm{var}(s)\cap\mathrm{bnd}(\varphi_1)=\varnothing\implies
        \frac{s}{y}[\varphi_1]\implies\frac{s}{y}[\varphi].
    \]

    \vspace{3pt}
    \item 假设$\varphi_1,\varphi_2$满足命题,
    再假设$\varphi=(\varphi_1\to\varphi_2)$满足
    $\mathrm{var}(s)\cap\mathrm{bnd}(\varphi)=\varnothing$. 那么
    由$\mathrm{bnd}(\varphi)=\mathrm{bnd}(\varphi_1)\cup\mathrm{bnd}(\varphi_2)$得
    \vspace{3pt}
    \[
        \mathrm{var}(s)\cap\mathrm{bnd}(\varphi_1)=\varnothing\land
        \mathrm{var}(s)\cap\mathrm{bnd}(\varphi_2)=\varnothing\implies
        \frac{s}{y}[\varphi_1]\land\frac{s}{y}[\varphi_2]\implies
        \frac{s}{y}[\varphi].
    \]

    \item 对任意的变元$x$, 假设$\varphi_1$满足命题,
    再假设$\varphi=\exists x\varphi_1$满足
    $\mathrm{var}(s)\cap\mathrm{bnd}(\varphi)=\varnothing$.
    
    \vspace{3pt}\hspace{2.17em}
    如果$y=x$, 那么$\displaystyle\frac{s}{y}[\varphi]$为真;

    \vspace{3pt}\hspace{2.17em}
    如果$y\neq x$, 首先由$\mathrm{var}(s)\cap\mathrm{bnd}(\varphi_1)=\varnothing$
    得$\displaystyle\frac{s}{y}[\varphi_1]$, 其次由$x\in\mathrm{bnd}(\varphi)$得
    $x\notin\mathrm{var}(s)$. 所以$\displaystyle\frac{s}{y}[\varphi]$.
    \hfill $\qedsymbol$
\end{enumerate}

% XZMH XDXH ~





\newpage

% 之前的版本是用单点替换的自由性来定义一般替换的自由性,
% 而现在的版本是直接定义一般的情况, 由此导出单点的情况.
% 因为这个定理, 所以这两种定义方式是等价的.
% 修改的时候觉得现在的版本更加贴近本质. 但是这样就要花篇幅证明这个定理.
% 更抽象地说, 第一版是给出两个定义, 而第二版是给出一个定义, 再证明一个定理.
% 现在想想, 似乎难以判断哪种方式更好.

% 相比第二版, 第一版不需要证明此定理,
% 但是需要额外证明"一般替换的自由性"满足它该有的性质的一个定理.
% 那个定理的证明似乎比这里的证明要简洁一些, 不太记得了.

% 第一版的概念感觉太冗杂, 会不得不写更多的简单但无聊的"定理",
% 并且不知道怎么给所谓"单点替换的自由性"一个好的名字.
% 抽象的(不是广泛使用的)名字大量出现会导致笔记的可读性很差. 所以最终决定保留第二版.
% 2026.02.28

\hypertarget{sec:定理1.5}{}
\noindent\textbf{定理1.5}

给定一阶语言$\mathcal{L}$. 对任意的公式$\varphi$和替换$\theta$,
$\theta[\varphi]$当且仅当$\displaystyle\forall y\in\mathcal{V}\cup\mathcal{C}:
\frac{\theta(y)}{y}[\varphi]$.

\noindent{\kaishu 证明.}

% 这些证明的书写真的很烦, 各种命题一层套一层, 不知道话该怎么说. 虽然结论很美妙.

\begin{enumerate}
    \item 对任意的原子公式$\varphi$, 两边始终为真.
    
    \item 假设$\varphi_1$满足命题, 则对$\varphi=\neg\varphi_1$有
    \[
        \theta[\varphi]\iff\theta[\varphi_1]\iff
        \forall y:\frac{\theta(y)}{y}[\varphi_1]\iff
        \forall y:\frac{\theta(y)}{y}[\varphi].
    \]

    \item 假设$\varphi_1,\varphi_2$满足命题, 则对$\varphi=(\varphi_1\to\varphi_2)$有
    \[
        \theta[\varphi]\iff\theta[\varphi_1]\land\theta[\varphi_2]\iff
        \forall y:\frac{\theta(y)}{y}[\varphi_1]\land\frac{\theta(y)}{y}[\varphi_2]
        \iff\forall y:\frac{\theta(y)}{y}[\varphi].
    \]

    \item 对任意的变元$x$, 假设$\varphi_1$满足命题,
    则对$\varphi=\exists x\varphi_1$有
    \vspace{3pt}
    \begin{equation*}\begin{aligned}
        \theta[\varphi]\iff&\theta/x[\varphi_1]\land
        \forall y\in\mathrm{at}(\varphi):x\notin\mathrm{var}\,\theta(y)\\
        \iff&\forall y:\frac{\theta/x(y)}{y}[\varphi_1]\land
        \forall y\in\mathrm{at}(\varphi):x\notin\mathrm{var}\,\theta(y)\\
        \iff&\forall y\neq x:\frac{\theta(y)}{y}[\varphi_1]\land
        \forall y\in\mathrm{at}(\varphi_1)\setminus\{x\}:
        x\notin\mathrm{var}\,\theta(y)\\
        \iff&\forall y\neq x:\left(\frac{\theta(y)}{y}[\varphi_1]\land
        (y\in\mathrm{at}(\varphi_1)\to x\notin\mathrm{var}\,\theta(y))\right)\\
        \iff&\forall y:\frac{\theta(y)}{y}[\varphi].
    \end{aligned}\tag*{\qedsymbol}\end{equation*}
\end{enumerate}

% 我收回前面的话. 这里的一长串的证明尝试了直接一路写过去, 而不是用自然语言讨论.
% 居然成功了. 很优雅.

\vspace{10pt}

使用这些定理可以得到替换自由的一个充分条件.

% 图穷匕见了. 写这一整节都是为了下面的这个定理, 用它证明替换定理的一个推论.
% (替换定理: 定理2.1)

% 值得注意的是, 黄书棋老师课程中的替换定理使用的就是
% &\forall y\in\mathrm{at}(\varphi):
% \mathrm{var}\,\theta(y)\cap\mathrm{bnd}(\varphi)=\varnothing&
% 这个条件, 它更容易判定. 但它只是真正的自由替换的一个充分条件.

% 【数理逻辑/模型论 2025 黄书棋老师 第十一节课：项与公式的替换；替换定理】
% https://www.bilibili.com/video/BV1fQSmBdEMr

\vspace{10pt}

\hypertarget{sec:定理1.6}{}
\noindent\textbf{定理1.6}

给定一阶语言$\mathcal{L}$. 对任意的公式$\varphi$和替换$\theta$, 如果
\[
    \forall y\in\mathrm{at}(\varphi):
    \mathrm{var}\,\theta(y)\cap\mathrm{bnd}(\varphi)=\varnothing,
\]
那么$\theta[\varphi]$.

\noindent{\kaishu 证明.}

对任意的变元或常元$y$:
\vspace{3pt}
\begin{enumerate}
    \item 如果$y\notin\mathrm{at}(\varphi)$,
    那么有$\displaystyle\frac{\theta(y)}{y}[\varphi]$,
    \item 如果$y\in\mathrm{at}(\varphi)$,
    依假设有$\mathrm{var}\,\theta(y)\cap\mathrm{bnd}(\varphi)=\varnothing$,
    从而也有$\displaystyle\frac{\theta(y)}{y}[\varphi]$.
\end{enumerate}

综上即得$\theta[\varphi]$. \hfill $\qedsymbol$

\end{document}